\documentclass[10pt]{article}
\usepackage[margin=20mm]{geometry}
\usepackage{amsmath,amssymb,mathtools,fontspec,unicode-math,xurl,hyperref}
\IfFontExistsTF{DejaVu Serif}{\setmainfont{DejaVu Serif}}{\setmainfont{Latin Modern Roman}}
\setmathfont{latinmodern-math.otf}
\setlength{\parindent}{0pt}\setlength{\parskip}{6pt}\sloppy
\title{Initial scientific review of
The singular sets of degenerate and nonlocal elliptic equations on Poincaré–Einstein manifolds}
\author{Scientific initial review}
\date{28 October 2026}
\begin{document}
\maketitle
Abstract: The review identifies an incorrect curvature coefficient in the extension equation, with a fully explicit counterexample on the hyperbolic four-ball. This is a Major issue because it invalidates the stated identification of the equation used in the curved argument with the fractional GJMS extension. Two further local formula errors are confirmed. A separate high-order extension step remains an evidence request, not an established counterexample to a theorem. The geometric singular-set conclusions themselves are not disproved by these checks.\par
\section*{Source, scope, and principal assessment}
This report concerns only Xumin Jiang, Yannick Sire, and Ruobing Zhang, arXiv:2309.09948v1, dated 18 September 2023, 47 PDF pages. The frozen project identifier is M25-0361; the associated journal DOI is https://doi.org/10.4171/jems/1740. No equivalence with the journal text is assumed. All page references below are one-based PDF pages, which agree with the printed numbering.\par
The initial review read the full source text. The present independent acceptance is confined to the reported Problems and their definitions, hypotheses, evidence, counterevidence, dependencies, and repair limits. Particular attention was given to the defining-function and scattering reduction, amplitude and scale dependence, the distinction between normalized and unnormalized frequencies, the cone-splitting and covering arguments, and the regularity required for the Hausdorff-measure step. These checks do not independently reprove every cited theorem or every compactness assertion. The detailed support inventory distinguishes textual reading, selected source-page inspection, exact symbolic computations, and remaining dependencies.\par
The paper works with \(n\geq2\), \(m=n+1\), \(\gamma\in(0,1)\), and \(a=1-2\gamma\in(-1,1)\). In the curved setting the compactification is specifically the Fefferman–Graham one. The spectral condition, the nonnegative Yamabe assumption in Section 4, the obstruction-flat qualification in even boundary dimension, and the distinction between smoothness up to one side and smoothness after doubling matter. In the flat section the authors require a Schwartz trace and a global energy extension. Local polynomial examples cannot silently replace those hypotheses.\par
Assessment: Major revision warranted for the exact reviewed version. Problem 1 establishes damage to the central extension identification, even though its coefficient has an explicit algebraic correction. Problems 2 and 3 are local Minor issues whose demonstrated effects are narrower. Problem 4 is classified Needs Evidence. The scientific decision remains pending and there are no human reviews. The independent acceptance checks recorded with this revised report concern the accuracy and scope of the existing assessment.\par
\section*{Problem 1: the curvature coefficient changes the extension equation}
Location and status. Lemma 2.11, pp. 11–12, especially (2.20), the display following (2.27), and (2.28); Proposition 2.13 and (2.31); the use of the same potential in (4.2) and (4.92). Confirmed; Major.\par
Equation reference: eq:printed. The starting expression (2.20) contains \(\gamma^2-\tfrac14\). The calculation on p. 12 instead uses \(\tfrac{\gamma^2}{4}-1\), and concludes \[
 C_{\mathrm{printed}}=\frac{n^2-3-4n\gamma+\gamma^2}{4n(n-1)}.\] This is not a harmless alternative normalization: the same \(a=1-2\gamma\) and the same scalar curvature are used on both pages.\par
\subsection*{Independent derivation}
Equation reference: eq:correct. Write \(\rho\) for the paper’s Fefferman–Graham defining function. Its stated identities give \[\Delta_{\bar g}\rho=-\frac{\rho R_{\bar g}}{n-1},
 \qquad
 1-|\nabla\rho|_{\bar g}^{2}=\frac{\rho^{2}R_{\bar g}}{n(n-1)}.\] Since \((a/2)(a/2-1)=\gamma^2-\tfrac14\), substituting into (2.20) yields \[\begin{aligned}
 E(\rho)
 &= -\frac a2\rho^{a-1}\Delta_{\bar g}\rho
   +\left(\gamma^2-\frac14\right)\rho^{a-2}
       (1-|\nabla\rho|_{\bar g}^{2})
   +\frac{n-1}{4n}R_{\bar g}\rho^a \notag\\
 &= C_{\mathrm{correct}}\rho^aR_{\bar g},\qquad
 C_{\mathrm{correct}}
 =\frac{a}{2(n-1)}
  +\frac{\gamma^2-\tfrac14}{n(n-1)}
  +\frac{n-1}{4n}
 =\frac{(n-2\gamma)^2}{4n(n-1)}.                 \end{aligned}\] Consequently \[C_{\mathrm{correct}}-C_{\mathrm{printed}}
   =\frac{3(\gamma^2+1)}{4n(n-1)}>0.\] The special case \(\gamma=\tfrac12\) is an immediate internal check: \(a=0\), so (2.20) itself reduces to \(E=(n-1)R_{\bar g}/(4n)\).\par
\subsection*{A complete admissible geometric example}
Let \(B^4\subset\mathbb R^4\), \(r=|x|\), \(n=3\), \(\gamma=\tfrac12\), and put \[g_+=\frac{4}{(1-r^2)^2}|dx|^2,\quad
 \phi=\frac{1-r^2}{2},\quad
 \rho=\frac{1-r^2}{2}e^\phi,\quad
 \bar g=e^{2\phi}|dx|^2.\] This is the complete hyperbolic four-ball with round conformal infinity \(S^3\), \(\operatorname{Ric}_{g_+}=-3g_+\). Direct radial differentiation gives \[-\Delta_{g_+}\log\rho=3,\qquad
 R_{\bar g}=6(4-r^2)e^{r^2-1}>0.\] Moreover, for the round geodesic defining function \(y=2(1-r)/(1+r)\) one has \(\rho/y\to1\) at the boundary. Thus this is the stipulated Fefferman–Graham compactification. The boundary is smooth and totally geodesic. The boundary Yamabe sign is positive and the hyperbolic spectral bottom is \(9/4>2=s(3-s)\), where \(s=2\). In this odd boundary dimension no obstruction-flat exception is needed.\par
Equation reference: eq:residual. Take the smooth boundary datum \(f\equiv1\) and \[u=\frac{1-r^2}{2},\qquad U=\rho^{s-n}u=e^{-\phi}.\] Then \(-\Delta_{g_+}u-2u=0\), \(U|_{S^3}=1\), and \(U\) is smooth on the compactified closed ball and belongs to the required energy space (\(a=0\)). No Schwartz assumption is involved in Proposition 2.13, which allows arbitrary smooth boundary data. The correct extension satisfies \[-\Delta_{\bar g}U+\frac16 R_{\bar g}U=0.\] In contrast, (2.28) gives \(C_{\mathrm{printed}}=1/96\), and the residual in the printed equation is \[
 -\Delta_{\bar g}U+\frac1{96}R_{\bar g}U
 =\frac{15}{16}(r^2-4)e^{\,3(r^2-1)/2}.\] At \(r=0\) this equals \(-\tfrac{15}{4}e^{-3/2}\), and is nonzero. This is an actual example under the extension-reduction hypotheses, not a formal unconstrained test function. The radial identities and residual were also independently checked with exact SymPy arithmetic; the retained script, output, and exit code are part of the evidence.\par
\subsection*{Impact, counterevidence, and repair boundary}
Proposition 2.13 incorporates the coefficient from Lemma 2.11 into the equation whose conormal derivative is identified with \(P_{2\gamma}f\). The explicit spectral extension above fails that printed equation. Later curved arguments repeatedly use \(J_{\bar g}=C_{n,\gamma}R_{\bar g}\), and the unique-continuation transfer in (4.86) again invokes the spectral interpretation. This establishes a failure of a central proof connection and warrants Major classification; it need not refute every singular-set theorem.\par
The flat model has \(R_{\bar g}=0\), so it does not reveal this error. Many estimates in Section 4 use only boundedness of a smooth lower-order potential, which is evidence that parts of the proof may survive the correction. Replacing the coefficient by [eq:correct] (\#eq:correct) is necessary. The extension correspondence, uniqueness/coercivity arguments, frequency definitions, and transfer to boundary solutions must then be checked with that corrected equation. The present review does not certify this entire repair or assert that the final dimension bounds are false.\par
\section*{Problem 2: the displayed Caccioppoli estimate has inconsistent homogeneity}
Location and status. Lemma 2.15, p. 14, (2.39), with the local equation (2.38) on p. 13. Confirmed local error; Minor.\par
For \(\beta>1\) the left side involves \(\nabla(\eta U^{\beta/2})\), whereas the right side contains \(U^2\) with a constant stated to depend only on \(r,a,\bar g\). Fix \(\beta=4\). In the Euclidean compactification of the hyperbolic upper-half-space model already described in the paper, take \(a=0\), \(J=0\), and the local energy solution \(U=A>0\). Fix a nonzero smooth cutoff supported strictly inside the unit ball and set \(I=\int|\nabla\eta|^2>0\). The asserted inequality would say \[A^4 I\leq C A^2 I.\] For any proposed uniform \(C\), choosing \(A^2>C\) contradicts it. All quantities are finite on the local ball. This is a counterexample to the local estimate, not to a global theorem requiring a Schwartz trace. More generally, multiplication of any nonzero admissible local solution by \(A\) produces the same incompatible amplitude powers whenever the cutoff expression is nonzero.\par
A direct corrected estimate can be obtained without assuming the disputed formula. For positive \(U\), \(\beta\geq2\), and \(w=\rho^a\), define \[A_0=\int w\eta^2 U^{\beta-2}|\nabla U|^2,\quad
 B=\int wU^\beta|\nabla\eta|^2,\quad
 D=\int w|J|\eta^2U^\beta.\] Testing the equation by \(\eta^2U^{\beta-1}\) and applying Young’s inequality gives \[A_0\leq\frac{4B}{(\beta-1)^2}+\frac{2D}{\beta-1},\] and hence \[\int w|\nabla(\eta U^{\beta/2})|^2
 \leq\left(2+\frac{2\beta^2}{(\beta-1)^2}\right)B
       +\frac{\beta^2}{\beta-1}D.\] For signed solutions the corresponding absolute-value regularization must be specified. This calculation supplies the correctly homogeneous input needed for the usual iteration. The displayed version cannot justify iteration as written, but this check has not shown a failure of the later regularity conclusion or the novel geometric estimates. That demonstrated scope, rather than ease of correction alone, is the reason for the Minor classification.\par
\section*{Problem 3: the sign of the radial weight derivative}
Location and status. Lemma 4.1, p. 32, (4.9)–(4.10). Confirmed; Minor.\par
For the full permitted interval \(a\in(-1,1)\) the paper displays \[\partial_r\rho^a=\frac{a\rho^a}{r}(1+O(r))\] and then asserts nonnegativity. Choose \(a=-\tfrac12\), equivalently \(\gamma=\tfrac34\). Along any inward ray on which \(\rho>0\), and for sufficiently small positive \(r\) such that \(|O(r)|<1/2\), the displayed derivative is strictly negative. In the flat normal-ray model \(\rho=r\) it is exactly \(-1/(2r^{3/2})\). This countercheck concerns the defining function and needs no unverified assumption about a solution \(U\).\par
The correct sign for small positive \(r\) is the sign of \(a\), with zero derivative for \(a=0\). The subsequent computations inspected on pp. 37–38 use the signed coefficient \(a+O(t)\) or \(a/r+O(rt)\), rather than the erroneous nonnegativity claim. This is counterevidence to a claim that the sign typo itself destroys the almost-monotonicity argument. It should be corrected, but this review does not establish such a downstream failure.\par
\section*{Problem 4: high-order extension across the doubled boundary}
Location and status. Proposition 4.20, pp. 43–45, especially (4.101)–(4.110), and its use in Theorem 4.19. Needs Evidence; not a confirmed theorem counterexample.\par
The negative-side formula in (4.110) reverses the signs of odd Taylor coefficients. For a nonzero odd coefficient \(c_\ell\) with \(\ell\leq2m^2\), the polynomial contribution to the two one-sided \(\ell\)th normal derivatives is respectively \[\ell!c_\ell\quad\hbox{and}\quad-\ell!c_\ell.\] A Taylor remainder with vanishing jets through that order cannot remove this discrepancy. Thus the statement that the displayed construction is clearly \(C^{2m^2}\) needs a proof that the relevant odd jets vanish, or a different extension construction with matching jets and controlled remainder. Smoothness on the positive side alone does not prove this. The finite regularity of the doubled coefficients, explicitly discussed in Remark 4.21, makes the distinction material.\par
Lemma 2.19 provides relevant counterevidence: (2.42)–(2.43), p. 14, state \(C^{k,\alpha}\) estimates for every \(k\), and the proof describes the estimates as holding across the doubled boundary. If these estimates apply with the stated two-sided regularity in Proposition 4.20, the evenness of \(U\) forces every odd normal derivative through order \(2m^2\) to vanish at the boundary. The conditional jet calculation would then present no obstruction to matching those derivatives.\par
The scope and justification of this regularity must therefore be reconciled with Remark 2.20, p. 15, which describes generically finite smoothness, and Remark 4.21, p. 43, which explains why the required high-order approximation cannot be obtained directly from the doubled coefficients. These statements do not by themselves establish that either regularity assertion or Proposition 4.20 is false. No admissible Poincaré–Einstein solution satisfying the frequency and approximation hypotheses with a nonzero offending jet has been constructed here. The assessment remains Needs Evidence, concerning the applicable regularity and the quantitative extension estimate, rather than a confirmed failure of the proposition or the Hausdorff bound.\par
There is also a quantifier discrepancy to clarify: Proposition 4.20 introduces \(\delta\), but (4.101) displays \(\epsilon\) in the approximation hypothesis, while (4.103) estimates by \(C_k\delta\). A revised proof should state the smallness parameter actually required. This is included in the same evidence request, not counted as an additional confirmed Major issue.\par
\section*{Coverage, limits, and conclusion}
The core path inspected runs from the spectral extension through weighted regularity, frequency control, quantitative rigidity and cone splitting, effective coverings, and the high-order stability step for Hausdorff measure. The source’s horizontal/mixed boundary decomposition was retained; codimension two for the horizontal part was not incorrectly applied to the entire trace singular set. The normalization and decay restrictions were checked before considering polynomial tests. Several tempting tests were rejected because they did not satisfy the global Schwartz/energy hypotheses.\par
The references on weighted regularity, quantitative stratification, and complex-analytic stability were identified through the paper, but their complete proofs were not independently re-established. In particular, the complex isolated-zero premise used in the stability reduction was not turned into a confirmed defect without the missing admissible evidence. The remark on rectifiability is likewise not independently certified here.\par
A bounded search of the exact arXiv record, an author research page, and queries for corrections, code, and supplements found no linked correction or separate implementation for this exact source during the check. This is a search limitation, not proof that no additional material exists. Acquisition metadata also records the inability to obtain the publisher PDF. No author code was executed; lack of code is not a mathematical defect.\par
The decisive result is the explicit failure of the curvature coefficient and of the extension identification using it. The repair boundary is stated above. Two narrower displayed formulas also require correction. The separate high-order extension question remains open in this review.\par
\section*{Supporting checks}
`4e62902d0df96aa76fa329aa0781ce7c64b9da01d2882eea550c081f3841057f`.\par
Backend equivalence was not independently verified. The initial review’s calculations, full-text reading, and page inspections retain their original attribution. In the present independent acceptance, the primary scientific conversation rechecked the reported Problems, corrected the discussion of counterevidence in Problem 4, and invoked the revised report’s compilation and actual page inspection. The primary verifier made the final judgments. Exact command and page records distinguish reused historical checks from new operations. Rendering alone is not counted as viewing; the scientific decision remains pending and the human-review list is empty.\par
9 X. Jiang, Y. Sire, and R. Zhang, The singular sets of degenerate and nonlocal elliptic equations on Poincaré–Einstein manifolds, arXiv:2309.09948v1 (18 September 2023), https://arxiv.org/abs/2309.09948v1. Exact reviewed source. arXiv, version record for 2309.09948, https://arxiv.org/abs/2309.09948, consulted 27 October 2026 for version and supplementary-material search. X. Jiang, research publication page, https://sites.google.com/site/xuminjiang60/research, consulted 27 October 2026 for related-material search.\par
\end{document}
